%==============================================================================
% Chapter 8 — Generative AI, NLP & Agentic Systems
% Curriculum: Phase 2, Chapter 8 of docs/AI_ML_Master_Checklist.md
% Companion code: projects/08_gen_ai/
%
% Section skeleton only. Figures for this chapter go in theory/figures/.
%==============================================================================
\chapter{Generative AI, NLP \& Agentic Systems}
\label{ch:gen-ai}

The longest chapter in the book, and the one that collects debts from everywhere
else.\projectref{projects/08\_gen\_ai/} Variational autoencoders are trained with the evidence lower bound of
\cref{sec:variational-inference}. Low-rank adaptation is the truncated SVD of
\cref{eq:eckart-young} applied to a weight update. Retrieval ranks documents by
the cosine similarity that fell out of \cref{eq:dot-product}, and perplexity is
the entropy of \cref{def:information} exponentiated.

The chapter runs from the foundational generative families through the
transformer and the techniques for adapting it, then out into the systems built
\emph{around} a frozen model --- retrieval, tool use, orchestration --- and
closes on how to tell whether any of it works, which is harder here than
anywhere else in the book.

%==============================================================================
\section{Variational Autoencoders}
\label{sec:vae}

\subsection{The Latent Variable Model}

\subsection{The Reparameterization Trick}

\subsection{The ELBO as a Training Objective}

%==============================================================================
\section{Generative Adversarial Networks}
\label{sec:gan}

\subsection{The Minimax Game}

\subsection{Training Instability and Mode Collapse}

%==============================================================================
\section{Diffusion Models}
\label{sec:diffusion}

\subsection{The Forward Noising Markov Chain}

\subsection{The Reverse Denoising Process}

\subsection{The Denoising U-Net}

\subsection{Classifier-Free Guidance}

%==============================================================================
\section{Text Representation}
\label{sec:text-representation}

\subsection{Tokenization and Subword Vocabularies}

\subsection{Embeddings}

%==============================================================================
\section{The Transformer Architecture}
\label{sec:transformer}

\subsection{Scaled Dot-Product Attention}

\subsection{Multi-Head Attention}

\subsection{Positional Encoding}

\subsection{The Full Block: Residuals, Normalisation and the Feed-Forward Layer}

\subsection{Encoder, Decoder and Decoder-Only Variants}

%==============================================================================
\section{Parameter-Efficient Fine-Tuning}
\label{sec:peft}

\subsection{Full Fine-Tuning and Why It Is Avoided}

\subsection{LoRA and Low-Rank Decomposition}

\subsection{QLoRA and Adapter Variants}

%==============================================================================
\section{Reinforcement Learning Fundamentals}
\label{sec:rl}

% Prerequisite for \cref{sec:alignment}, and the one paradigm in the book that
% does not fit the supervised template of \cref{ch:deep-learning}: the model
% generates its own data and the objective is a return rather than a label.
% PPO is a variance fix to a policy gradient; without the chain below it is an
% algorithm you can invoke but not reason about.

\subsection{Markov Decision Processes}

\subsection{Returns, Discounting and the Objective}

\subsection{Value and Action-Value Functions}

\subsection{The Bellman Equations}

\subsection{Dynamic Programming: Policy and Value Iteration}

\subsection{Temporal-Difference Learning and Q-Learning}

\subsection{Deep Q-Networks}

\subsection{The Policy Gradient Theorem and REINFORCE}

\subsection{Baselines and Variance Reduction}

\subsection{Actor-Critic and Advantage Estimation}

\subsection{Trust Regions: from TRPO to the Clipped Surrogate}

%==============================================================================
\section{Alignment}
\label{sec:alignment}

\subsection{Supervised Fine-Tuning}

\subsection{Reward Modeling from Preference Pairs}

\subsection{RLHF with PPO}

% Builds directly on \cref{sec:rl}: the policy is the language model, the action
% space is the vocabulary, and the reward comes from the model fitted above ---
% plus a KL penalty against the reference policy, which is
% \cref{def:information} doing work again.

\subsection{Direct Preference Optimization}

% The derivation that removes the RL loop entirely by solving for the optimal
% policy in closed form and substituting it into the reward model's likelihood.

%==============================================================================
\section{Local Inference and Quantization}
\label{sec:quantization}

\subsection{Numeric Formats from FP32 to INT4}

\subsection{Post-Training Quantization and Quantization-Aware Training}

\subsection{KV Cache Mechanics}

%==============================================================================
\section{The Hugging Face Ecosystem}
\label{sec:huggingface}

\subsection{\texttt{transformers} and \texttt{datasets}}

\subsection{Hub Management and Model Cards}

%==============================================================================
\section{Prompt Engineering}
\label{sec:prompting}

\subsection{Zero-Shot and Few-Shot Prompting}

\subsection{Chain-of-Thought}

\subsection{Tree-of-Thought}

%==============================================================================
\section{Retrieval-Augmented Generation}
\label{sec:rag}

\subsection{Chunking Strategies}

\subsection{Embedding Models}

\subsection{Vector Databases}

\subsection{Distance Metrics}

\subsection{The Ingestion-to-Answer Pipeline}

\subsection{Reranking and Hybrid Search}

%==============================================================================
\section{Agentic Systems}
\label{sec:agents}

\subsection{The ReAct Framework}

\subsection{Tool and Function Calling}

\subsection{Multi-Agent Orchestration as a DAG}

%==============================================================================
\section{Generative Evaluation and Benchmarking}
\label{sec:gen-eval}

\subsection{N-Gram Metrics: BLEU and ROUGE}

\subsection{Perplexity and Cross-Entropy}

\subsection{RAG Metrics: Context Precision, Recall and Faithfulness}

\subsection{LLM-as-Judge and Its Failure Modes}

%==============================================================================
\section{Adversarial Robustness and AI Security}
\label{sec:ai-security}

% \Cref{ch:capstone} puts a model behind a public HTTP endpoint, which turns
% everything below from a research topic into an operational requirement.

\subsection{Adversarial Examples: FGSM and PGD}

% The gradient of the loss with respect to the \emph{input} rather than the
% weights --- the same machinery as \cref{sec:pinn}, pointed at an attack.

\subsection{Data Poisoning and Backdoor Attacks}

\subsection{Prompt Injection and Jailbreaks}

\subsection{Indirect Injection Through Retrieved Documents}

% The failure mode specific to \cref{sec:rag}: retrieved context is untrusted
% input, but the model reads it as instructions.

\subsection{Model Extraction and Membership Inference}

\subsection{PII Handling, Data Minimisation and Differential Privacy}

\subsection{Input Validation and Output Guardrails}
